%%%%%%%%%%%%%%%%%%%%%%%%%%%%%%%%%%%%%%%%%%%%%%%%%%%%%%%%%%%%
%% CHAPTER 14 -- What to Ask Before You Deploy
%%%%%%%%%%%%%%%%%%%%%%%%%%%%%%%%%%%%%%%%%%%%%%%%%%%%%%%%%%%%

\chapter{What to Ask Before You Deploy}
\label{ntg:ch14}

\section*{A note on purpose}
\addcontentsline{toc}{section}{A note on purpose}

The preface promised that a reader who finished this guide
would be able to ask better questions of the engineers who
build these systems, the vendors who sell them, and the
colleagues who deploy them. The thirteen chapters that
came before have built the vocabulary and the mechanisms.
This chapter is the bridge from understanding to use.

It is organised around the failure modes the guide has
already covered. For each, there are two or three questions
a non-technical reader can ask that are, in fact, demanding
questions --- ones that require a real answer and that reveal
something about the quality of the system. They are not
trick questions. They are the questions a competent and
informed person asks when deciding whether to trust a system
enough to use it in a context where the answer matters.

A vendor, engineer, or colleague who cannot answer these
questions --- or who dismisses them as too technical, too
obscure, or not relevant to your use case --- is telling
you something worth knowing.

\section{On hallucination}
\addcontentsline{toc}{section}{On hallucination}

A language model produces the most plausible continuation
of the prompt it has been given. Plausibility and truth are
correlated but not the same. The model has no native way to
tell the two apart. The grounding axis of the error model
in Chapter~13 is where hallucinations live.

\medskip
\noindent\textbf{Ask:}

\begin{itemize}
  \item \emph{What is the retrieval architecture, and how
        are retrieved passages attributed in the response?}
        A system that generates responses from the model's
        weights alone is systematically more prone to
        hallucination than one that retrieves from a
        specific document store and cites its sources.
        If citations are present, you can verify them.
        If they are not, you cannot.

  \item \emph{How was the system evaluated for factual
        accuracy on content in our specific domain, and
        what was the error rate?} General benchmarks
        measure performance on broadly sampled questions.
        Your deployment is not broadly sampled; it is
        concentrated in your domain. A system with
        impressive general-purpose accuracy may perform
        far worse on specialised content that was
        underrepresented in training. Ask for domain-specific
        numbers.

  \item \emph{What happens when the model is asked about
        something it does not know?} A well-designed system
        should express uncertainty and decline to confabulate
        when it lacks grounding. A poorly calibrated system
        will confabulate confidently. Ask for examples of
        the system being asked questions it should not be
        able to answer. The behaviour in those cases tells
        you more than the behaviour in the easy cases.
\end{itemize}

\section{On sycophancy and reward hacking}
\addcontentsline{toc}{section}{On sycophancy and reward hacking}

The alignment pipeline optimises the model against a reward
model trained on human preferences. If the preference data
rewarded agreement with the user, the policy will learn to
agree. If the preference data rewarded long responses, hedged
prose, or certain opening formulae, the policy will learn
those patterns too. These are features of the training
process, not bugs the model is hiding.

\medskip
\noindent\textbf{Ask:}

\begin{itemize}
  \item \emph{Has the model been tested for sycophancy
        specifically --- for example, by presenting it with
        false premises and measuring how often it agrees?}
        Sycophancy testing is not standard in most vendor
        evaluation suites. If a vendor has done it and has
        results to share, that is a mark of seriousness.
        If the question draws a blank, the system has
        probably not been evaluated for this failure mode.

  \item \emph{Does the model change its answer when the
        user pushes back, and if so, under what conditions?}
        A model that revises a correct answer in response
        to user disagreement is behaving sycophantically.
        A model that revises an incorrect answer when
        presented with new evidence is behaving correctly.
        The distinction is easy to miss and easy to test.
        Ask for the policy and ask to see examples.

  \item \emph{What KL constraint was used during RLHF
        training, and how was the trade-off between
        helpfulness and capability degradation monitored?}
        This question does not require the vendor to answer
        in detail. It requires them to acknowledge that
        the trade-off exists and that they tracked it.
        A vendor who has not thought about this has
        produced an aligned model without understanding
        what the alignment cost.
\end{itemize}

\section{On prompt brittleness}
\addcontentsline{toc}{section}{On prompt brittleness}

Two prompts that a human would read as equivalent can
produce qualitatively different outputs from the same model.
This is not random noise. It is a structural property of
the attention mechanism (Chapter~5): the model is sensitive
to the exact sequence of tokens it receives, and small
changes to that sequence change the weighting across the
context window.

\medskip
\noindent\textbf{Ask:}

\begin{itemize}
  \item \emph{Has the system been tested with paraphrased
        variants of the prompts used in production, and
        how much does output quality vary across those
        variants?} A brittle system will show high variance.
        A robust system will show low variance. Most
        evaluation suites test a single prompt formulation
        per task. Paraphrase testing is the exception.

  \item \emph{If the system uses a fixed system prompt,
        how was that prompt selected, and how often is
        it reviewed?} The system prompt encodes the
        situational context the model receives before the
        user's input. A poorly written system prompt can
        produce systematic failures in every interaction.
        Ask to see it. Ask how it was arrived at. Ask
        what happens when a user's input conflicts with it.

  \item \emph{What is the expected behaviour when users
        phrase questions in ways the system prompt did
        not anticipate?} Most real deployments receive
        inputs the designers did not imagine. The
        question is not whether this will happen but
        whether the system degrades gracefully or fails
        in unexpected ways.
\end{itemize}

\section{On context loss and long documents}
\addcontentsline{toc}{section}{On context loss and long documents}

A model with a large context window can accept a large
document as input. Accepting a document and reliably using
all of it are different things. The attention mechanism
described in Chapter~5 is theoretically capable of relating
any position in the context to any other, but in practice
the model's ability to use information from the middle of
a very long context is weaker than its ability to use
information from the beginning or end. This is the
``lost in the middle'' failure documented in Chapter~10.

\medskip
\noindent\textbf{Ask:}

\begin{itemize}
  \item \emph{How was the model evaluated on long-context
        tasks, and what was the tested context length?}
        A vendor may advertise a 128,000-token context
        window and evaluate the system only on tasks that
        fit in 8,000 tokens. The advertised window and the
        reliable window are not the same.

  \item \emph{For document-based tasks, where is the
        relevant information typically placed in the
        prompt?} There is evidence that placing critical
        information at the beginning or end of a long
        prompt improves retrieval. If the system places
        it in the middle and has not been tested on this,
        the design may be suboptimal in a way the
        vendor has not measured.

  \item \emph{What happens when the context window is
        exceeded?} Some systems truncate silently.
        Others summarise. Others return an error. Truncation
        that the user does not see is a context failure
        that is invisible by design. Ask what the policy is
        and whether the user is notified.
\end{itemize}

\section{On prompt injection}
\addcontentsline{toc}{section}{On prompt injection}

A prompt injection attack occurs when a malicious party
embeds instructions in content the model processes ---
a retrieved document, a tool result, a user-supplied input
--- that cause the model to execute those instructions as
if they were part of the system prompt. It is, as
Chapter~12 notes, roughly what SQL injection once was to
web applications: a class of attack that arises from
treating untrusted input as trusted instruction.

\medskip
\noindent\textbf{Ask:}

\begin{itemize}
  \item \emph{Does the system process content from untrusted
        sources --- web pages, user documents, email,
        third-party tool outputs --- and if so, what
        mitigations are in place against prompt injection?}
        A system that retrieves from the public web and
        then generates responses from retrieved content
        is exposed to any attacker who can get a document
        indexed. Ask what the system does when a retrieved
        document contains instructions addressed to the
        model.

  \item \emph{What is the trust boundary between the
        system prompt and user input, and is it enforced
        at the model level or only at the application level?}
        Application-level enforcement (filtering, sanitising
        inputs before they reach the model) is weaker
        than model-level enforcement (the model having been
        trained to distinguish between instruction sources).
        Neither is perfect. Ask which is being relied on.

  \item \emph{Has the system been red-teamed against
        prompt injection specifically, and what was the
        result?} Red-teaming is the practice of deliberately
        trying to make the system behave badly. A vendor
        who has red-teamed and has results to share is
        ahead of one who has not. A vendor who is unaware
        of prompt injection as a threat class is, for most
        production deployments, not a suitable partner.
\end{itemize}

\section{On automation complacency}
\addcontentsline{toc}{section}{On automation complacency}

As Chapter~11 describes, the most insidious failure mode of
any reliable-seeming system is the erosion of the human
review that was supposed to catch its errors. A system that
is correct ninety-five per cent of the time is one that is
wrong one time in twenty, and a reviewer who has learned to
expect correctness will stop reading carefully. The failure
rate is a constant; the detection rate falls.

\medskip
\noindent\textbf{Ask:}

\begin{itemize}
  \item \emph{What human review process is in place for
        the outputs this system generates, and how has that
        process been designed to remain effective as the
        system is perceived as reliable?} The question is
        not only whether there is a review process but
        whether the process has been designed to resist
        complacency. Review processes that do not build
        in adversarial checks --- deliberately introducing
        known errors for reviewers to catch --- tend to
        degrade over time.

  \item \emph{How are errors discovered and fed back into
        the system or its operating procedures?} An
        isolated error in a single output is an incident.
        A pattern of errors is a systemic failure. The
        difference between organisations that learn from
        deployment and organisations that do not is whether
        there is a structured path from observed error to
        corrective action. Ask what that path is.

  \item \emph{Who is accountable for errors that pass
        through human review?} The question is not only
        about process but about incentives. If accountability
        for errors is diffuse, review will be less careful
        than if a specific person or team is responsible
        for what passes. This is an organisational question
        as much as a technical one.
\end{itemize}

\section{A closing note on these questions}
\addcontentsline{toc}{section}{A closing note on these questions}

The questions above are not a compliance checklist. A vendor
who has answered all of them satisfactorily has not thereby
produced a safe system; they have produced a system whose
known failure modes have been attended to. New failure modes
will emerge in deployment. The questions are the beginning
of a productive ongoing conversation, not the end of it.

A non-technical reader who can ask these questions, and who
knows enough to evaluate the answers, has something more
valuable than a checklist: a working model of how these
systems fail. That model is what the preceding chapters were
built to provide. The questions are its application.

The next and final chapter of this guide returns to the
reader who began it.
